\documentclass[11pt,a4paper]{article}
\usepackage[italian]{babel}
\usepackage[margin=2.5cm]{geometry}
\usepackage{parskip}
\usepackage{amsmath, amssymb, amsthm}
\usepackage{booktabs}
\usepackage{array}
\usepackage{xcolor}
\usepackage{needspace}
\usepackage{enumitem}
\usepackage{listings}
\usepackage{tikz}
\usepackage[most]{tcolorbox}
\usepackage[hidelinks]{hyperref}
\usetikzlibrary{decorations.pathreplacing, shapes.geometric, arrows.meta, positioning}

% Niente vedove/orfane: i paragrafi non lasciano righe isolate a fine/inizio pagina
\widowpenalty=10000
\clubpenalty=10000

% Elenchi più compatti
\setlist{itemsep=2pt, parsep=0pt, topsep=4pt, partopsep=0pt}

% --- Riquadri con bordo nero, senza numero (si spezzano tra due pagine) ---
% Uso: \begin{definizione} ... \end{definizione}
%      \begin{teorema}[Nome] ... \end{teorema}  (il nome è facoltativo)
\tcbset{nota/.style={enhanced jigsaw, breakable, colback=white,
  colframe=black, boxrule=0.5pt, sharp corners}}
\NewTColorBox{definizione}{}{nota,
  before upper={\textbf{Definizione.}\ }}
\NewTColorBox{teorema}{o}{nota,
  before upper={\textbf{Teorema\IfValueT{#1}{ (#1)}.}\ }}
\NewTColorBox{proposizione}{o}{nota,
  before upper={\textbf{Proposizione\IfValueT{#1}{ (#1)}.}\ }}
\NewTColorBox{proprieta}{o}{nota,
  before upper={\textbf{Proprietà\IfValueT{#1}{ (#1)}.}\ }}
\NewTColorBox{assioma}{o}{nota, boxrule=1pt,
  before upper={\textbf{Assioma\IfValueT{#1}{ (#1)}.}\ }}

% --- Ambienti semplici (senza riquadro) ---
% Il titolo sta in cima al blocco, anche se il contenuto inizia con un riquadro o
% con codice; e non resta mai da solo in fondo a una pagina.
% Uso: \begin{esempio}[Titolo facoltativo] ... \end{esempio}
\newcommand{\newrunin}[2]{%
  \NewDocumentEnvironment{#1}{o}{%
    \par\Needspace{4\baselineskip}\addvspace{\medskipamount}\noindent
    \textbf{#2}\IfValueT{##1}{ (##1)}\textbf{.}\ \ignorespaces}%
  {\par\addvspace{\medskipamount}}}
\newrunin{esempio}{Esempio}
\newrunin{esercizio}{Esercizio}
\NewTColorBox{osservazione}{}{nota,
  before upper={\textbf{Osservazione.}\ }}

% V e F nelle tavole di verità
\newcommand{\V}{\textsf{V}}
\newcommand{\F}{\textsf{F}}

% --- Codice C ---
% Uso: \begin{codice} ... \end{codice}      codice C numerato
%      \begin{comando} ... \end{comando}    comandi da terminale
%      \begin{risultato} ... \end{risultato} output di un programma
%      \lstinline|printf("ciao");|           codice dentro al testo
\lstdefinestyle{c}{
  language=C,
  basicstyle=\ttfamily\small,
  keywordstyle=\color{blue}\bfseries,
  commentstyle=\color{gray}\itshape,
  stringstyle=\color{red!70!black},
  numbers=left,
  numberstyle=\tiny\color{gray},
  numbersep=8pt,
  columns=flexible,
  keepspaces=true,
  breaklines=true,
  showstringspaces=false,
  tabsize=4,
  morekeywords={include, define},
  % lettere accentate nei commenti
  literate={à}{{\`a}}1 {è}{{\`e}}1 {é}{{\'e}}1 {ì}{{\`i}}1
           {ò}{{\`o}}1 {ù}{{\`u}}1 {È}{{\`E}}1 {À}{{\`A}}1
}
\lstset{style=c}
\newtcblisting{codice}{listing only, nota, left=6mm,
  listing options={style=c}}
\newtcblisting{comando}{listing only, nota,
  listing options={basicstyle=\ttfamily\small, numbers=none, language={},
    columns=flexible, keepspaces=true, breaklines=true}}
\newtcblisting{risultato}{listing only, enhanced jigsaw, breakable,
  colback=black!5, colframe=black, boxrule=0.5pt, sharp corners,
  listing options={basicstyle=\ttfamily\small, numbers=none, language={},
    columns=flexible, keepspaces=true, breaklines=true}}

% --- Diagrammi di flusso (simboli standard) ---
\tikzset{
  terminale/.style = {ellipse, draw, minimum width=2.5cm, minimum height=0.9cm},
  io/.style        = {trapezium, trapezium left angle=70, trapezium right angle=110,
                      draw, minimum width=2.5cm, minimum height=0.9cm},
  processo/.style  = {rectangle, draw, minimum width=2.5cm, minimum height=0.9cm},
  decisione/.style = {diamond, draw, aspect=2, minimum width=2.5cm},
  freccia/.style   = {thick, -{Stealth}}
}

\title{Ottimizzazione: mappe di Karnaugh}
\author{Marco Cerbella}
\date{Lezione 5}

\begin{document}
\maketitle

\section{Che cos'è una mappa di Karnaugh}

Le mappe di Karnaugh sono un metodo \emph{grafico} per semplificare le espressioni booleane. Furono inventate nel 1953 da Maurice Karnaugh, ingegnere delle telecomunicazioni dei Bell Labs, e sono molto efficaci per ottimizzare reti combinatorie con poche variabili (fino a quattro).

\begin{definizione}[Mappa di Karnaugh]
Una \emph{mappa di Karnaugh} è una rappresentazione formale, alternativa alla tabella di verità, di una funzione booleana: una griglia con una cella per ogni configurazione delle variabili, in cui si scrive il valore della funzione.
\end{definizione}

\begin{esempio}
Sia $f(x,y,z)$ con $f(001) = f(011) = 1$ e $f = 0$ altrove. Le righe della mappa indicano $x$ e le colonne indicano $yz$:
\begin{center}
\begin{tabular}{c|cccc}
 & $00$ & $01$ & $11$ & $10$ \\
\midrule
$0$ & 0 & 1 & 1 & 0 \\
$1$ & 0 & 0 & 0 & 0 \\
\end{tabular}
\end{center}
Le due celle con $1$ sono quelle di $f(001)$ e $f(011)$.
\end{esempio}

\section{Codice di Gray}

Nella mappa le configurazioni delle variabili non seguono l'ordine binario naturale ($00, 01, 10, 11$) ma un \emph{codice di Gray} (o codice a specchio): $00, 01, 11, 10$. Passando da una colonna alla successiva cambia un solo bit. Ciò è fondamentale perché mette in celle \emph{adiacenti} (con un lato in comune) i mintermini che danno luogo a semplificazioni.

Per esempio, nella mappa precedente
\[
x'y'z + x'yz = x'z\,(y' + y) = x'z.
\]

\section{Perché funzionano}

Due mintermini sono \emph{logicamente adiacenti} se sono uguali a meno di una variabile, che compare in forma naturale in uno e negata nell'altro (per esempio $x_1x_2'x_3$ e $x_1x_2x_3$). Se $p_1 + p_2$ è la somma di due mintermini logicamente adiacenti, allora $p_1 + p_2$ si può semplificare:
\[
x_1x_2'x_3 + x_1x_2x_3 = x_1x_3\,(x_2' + x_2) = x_1x_3.
\]
In una mappa di Karnaugh, a celle geometricamente adiacenti corrispondono mintermini logicamente adiacenti.

\section{Come si usa}

\begin{enumerate}
    \item Si raggruppano gli $1$ che compaiono in celle adiacenti: ogni gruppo è un \emph{$k$-cubo}.
    \item Da ogni $k$-cubo si ricava un \emph{implicante}, cioè un termine prodotto $p$ tale che, se $p$ vale $1$, allora $f$ vale $1$. L'implicante è il prodotto delle variabili che compaiono sempre in forma naturale (oppure sempre in forma negata) in tutte le celle del $k$-cubo.
    \item Coperti tutti gli $1$, si sommano gli implicanti: si ottiene un'espressione ottimizzata, in forma di somma di prodotti, corrispondente a $f$.
\end{enumerate}

\subsection*{Regole per generare i $k$-cubi}

\begin{enumerate}
    \item La mappa si immagina \emph{arrotolata} su se stessa: l'ultima colonna è adiacente alla prima e l'ultima riga è adiacente alla prima.
    \item I lati dei rettangoli in cui si raccolgono gli $1$ devono essere potenze di $2$.
    \item Ogni $k$-cubo deve essere il più grande possibile.
    \item Ogni $1$ deve stare in almeno un $k$-cubo.
    \item Un $1$ può stare in più $k$-cubi, se serve per coprire tutti gli $1$.
\end{enumerate}

\section{Esempi}

\begin{esempio}[Tre variabili]
Mappa con righe $x$ e colonne $yz$ e $f = y'z' + yz$:
\begin{center}
\begin{tabular}{c|cccc}
 & $00$ & $01$ & $11$ & $10$ \\
\midrule
$0$ & 1 & 0 & 1 & 0 \\
$1$ & 1 & 0 & 1 & 0 \\
\end{tabular}
\end{center}
I due $k$-cubi sono la colonna $00$ (implicante $y'z'$) e la colonna $11$ (implicante $yz$).
\end{esempio}

\begin{esempio}[Tre variabili, con la colonna $10$]
\begin{center}
\begin{tabular}{c|cccc}
 & $00$ & $01$ & $11$ & $10$ \\
\midrule
$0$ & 0 & 1 & 1 & 1 \\
$1$ & 1 & 1 & 0 & 1 \\
\end{tabular}
\end{center}
Si hanno tre $k$-cubi: le celle $01,11$ della riga $0$ (implicante $x'z$), le celle $00,01$ della riga $1$ (implicante $xy'$) e la colonna $10$ (implicante $yz'$). Quindi
\[
f(x,y,z) = x'z + xy' + yz'.
\]
\end{esempio}

\begin{esempio}[Quattro variabili]
Righe $xy$, colonne $zw$:
\begin{center}
\begin{tabular}{c|cccc}
 & $00$ & $01$ & $11$ & $10$ \\
\midrule
$00$ & 0 & 0 & 0 & 1 \\
$01$ & 0 & 1 & 0 & 1 \\
$11$ & 1 & 1 & 0 & 0 \\
$10$ & 1 & 1 & 0 & 0 \\
\end{tabular}
\end{center}
I $k$-cubi sono: le righe $11$ e $10$ nelle colonne $00,01$ (implicante $xz'$); le celle $01$ e $11$ della colonna $01$ (implicante $yz'w$); le celle $00$ e $01$ della colonna $10$ (implicante $x'zw'$). Quindi
\[
f(x,y,z,w) = xz' + yz'w + x'zw'.
\]
\end{esempio}

\begin{esempio}[Quattro variabili, con celle ai bordi]
\begin{center}
\begin{tabular}{c|cccc}
 & $00$ & $01$ & $11$ & $10$ \\
\midrule
$00$ & 1 & 0 & 0 & 1 \\
$01$ & 0 & 1 & 0 & 1 \\
$11$ & 1 & 1 & 0 & 0 \\
$10$ & 1 & 1 & 0 & 0 \\
\end{tabular}
\end{center}
Rispetto all'esempio precedente la cella $(00, 00)$ vale $1$. Va raccolta con la cella $(10, 00)$, perché la prima e l'ultima riga sono adiacenti: si ottiene l'implicante $y'z'w'$. Quindi
\[
f(x,y,z,w) = xz' + yz'w + x'zw' + y'z'w'.
\]
\end{esempio}

\begin{esempio}[Quattro variabili, cubo grande]
\begin{center}
\begin{tabular}{c|cccc}
 & $00$ & $01$ & $11$ & $10$ \\
\midrule
$00$ & 1 & 0 & 1 & 1 \\
$01$ & 0 & 0 & 1 & 1 \\
$11$ & 0 & 0 & 1 & 1 \\
$10$ & 1 & 0 & 1 & 1 \\
\end{tabular}
\end{center}
Le colonne $11$ e $10$ formano un $k$-cubo $4 \times 2$ (implicante $z$). Le quattro celle d'angolo (righe $00,10$ e colonne $00,10$) formano un secondo $k$-cubo, perché la mappa è arrotolata sia in orizzontale sia in verticale (implicante $y'w'$). Quindi
\[
f(x,y,z,w) = z + y'w'.
\]
\end{esempio}

\end{document}
